\subsection{Local limit theorems}
\label{ss:llt}

We now establish two local limit estimates concerning the random walks $\Sbn$ and $\Scn$
(Lemmas~\ref{lem:ll1} and \ref{lem:ll2}).



\paragraph*{Local limit theorem for $\Sbn$.}


Let $\phibn(t)=\Esb{e^{it\Sbn_{1}}}= \tfrac{F( \bbn e^{it})}{F(\bbn)}$  be the characteristic function of $\Sbn_{1}$,
i.e.\ of the distribution $\mubn$.
(Recall that $F(z)=\sum_{k \ge 0} \tfrac{(k+1)^{k-1}}{k!} z^k$ was introduced in the proof of Lemma~\ref{lem:exist}.)

\begin{lemma}\label{lem:phib}
Assume that $\tfrac{K_{n}}{n} \rightarrow 0$ and $K_{n} \rightarrow \infty$ as $n \rightarrow \infty$. For every $n$ sufficiently large, for every $ 0 \leq |t| \leq \pi$  we have
\begin{equation}
\label{eq:phib}  \ln | \phibn (t)| \leq -  \frac{K_{n}}{ 8n} t^{2}.
\end{equation}
\end{lemma}

\begin{proof}
We take two different approaches depending on the distance of $|t|$ to $0$.

First assume that $|t| \leq 1$. We have (see e.g.~\cite[Lemma 3.3.7]{Dur10})
\[\phibn(t)=1+\textrm{i}\E[\Sbn_{1}]t- \E[(\Sbn_{1})^{2}] \frac{t^{2}}{2}+r^{n}_{1}(t),\]
with $|r^{n}_{1}(t)| \leq |t|^{3} \E[(\Sbn_{1})^{3}]/6$ for every $n \geq 1$ and $t \in \R$. By using the expansion of $F$ around $0$, we see that $ \E[\Sbn_{1}] \sim \E[(\Sbn_{1})^{2}] \sim \E[(\Sbn_{1})^{3}] \sim \frac{K_{n}}{n} \rightarrow 0$ as $n \rightarrow \infty$, so that
\begin{align*}
  \hspace{ -2mm} |\phibn(t)|^{2} &= \left( 1+\textrm{i}\E[\Sbn_{1}]t- \E[(\Sbn_{1})^{2}] \frac{t^{2}}{2}+r^{n}_{1}(t) \right) \left( 1-\textrm{i}\E[\Sbn_{1}]t- \E[(\Sbn_{1})^{2}] \frac{t^{2}}{2}+\overline{r^{n}_{1}(t)} \right)\\
&= 1 - (\sbn)^{2} t^{2}+r^{n}_{2}(t),
\end{align*}
where $|r^{n}_{2}(t)| \leq  |t|^{3} K_{n}/(2n)$ for every $n$ sufficiently large, uniformly for $t \in [-1,1]$.
Therefore
\[ 2 \ln |\phibn(t)| \leq - (\sbn)^{2} t^{2}+ |t|^{3} K_{n}/(2n)= - \frac{K_{n}}{n} t^{2} \left(  \frac{(\sbn)^{2} n}{K_{n}}-|t|/2 \right)
\]
for every $n$ sufficiently large, uniformly for $t \in [-1,1]$.
Since $(\sbn)^{2} \sim K_{n}/n$ as $n \rightarrow \infty$ by Lemma~\ref{lem:exist} (ii), it follows that for every $n$ sufficiently large and $0 \leq |t| \leq 1$ we have $ \ln | \phibn (t)| \leq -  \frac{K_{n}}{ 8n} t^{2}$.

Now assume that $1 \leq |t| \leq \pi$. Since $F(z)=1+z+o(z)$ as $z \rightarrow 0$ and, by definition, $\phibn(t)=\tfrac{F( \bbn e^{it})}{F(\bbn)}$,
we have
\[|\phibn(t)|^{2}=(1+\bbn(e^{it}-1)+o(\bbn))(1+\bbn(e^{-it}-1)+o(\bbn))=1+2\bbn(\cos(t)-1)+o(\bbn),\]
where the $o(\bbn)$ is uniform in $1\leq |t| \leq \pi$.
Therefore, for every $1 \leq |t| \leq \pi$, for every $n$ sufficiently large,
\[ 2 \ln |\phibn(t)| \leq - \bbn \frac{t^{2}}{3}+o(\bbn) \leq - \frac{K_{n}}{4n} t^{2},\]
where for the last inequality we have used the fact that $\bbn \sim \frac{K_{n}}{n}$ as $n \rightarrow \infty$ by Lemma~\ref{lem:exist} (ii) and that $t$ is not too close to $0$.
\end{proof}





\begin{proof}[Proof of Lemma~\ref{lem:ll1}]
We first check that the convergence
\begin{equation}
\label{eq:cv1b}\E\bigg[e^{i t \frac{\Sbn_{N}-N \frac{K_{n}+1}{n-K_{n}} }{ \DbnN }} \bigg]  \quad \mathop{\longrightarrow}_{n \rightarrow \infty} \quad e^{- \frac{t^{2}}{2}}
\end{equation}
holds uniformly for $un \leq N \leq n$ and $t$ in compact subsets of $\R$.  As in the proof of Lemma~\ref{lem:phib}, since $\tfrac{\E[(\Sbn_{1})^{3}]}{(\DbnN)^{3}}=\mathcal O(\tfrac{1}{n \sqrt{K_{n}}})= o(\frac{1}{n})$, we have
\[\phibn\left( \frac{t}{ \DbnN } \right)=1+\textrm{i}\E[\Sbn_{1}] \frac{t}{ \DbnN } - \E[(\Sbn_{1})^{2}] \frac{t^{2}}{2(\DbnN)^{2}}+ o \left( \frac{1}{n} \right),\]
where the $o$ is uniform when $un \leq N \leq n$ and $t$ belongs to a compact subset of $\R$, so that
\[\ln \phibn \left( \frac{t}{ \DbnN } \right)=it  \frac{K_{n}+1}{(n-K_{n})\DbnN} - \frac{t^{2}}{2} \cdot \frac{(\sbn)^{2}}{(\DbnN )^{2}}+ o \left( \frac{1}{n} \right).\]
Thus
\[\Es{e^{i t \frac{\Sbn_{N}-N \frac{K_{n}+1}{n-K_{n}} }{ \DbnN }}}= \exp \left(- \frac{t^{2}}{2} \cdot  \frac{ N (\sbn)^{2} }{(\DbnN)^{2}} + o \left( \frac{N}{n} \right)   \right).\]
Since $(\DbnN)^{2}= N (\sbn)^{2}$, this  implies \eqref{eq:cv1b}.


We then follow the steps of the analytic proof of the standard local limit theorem for a sequence of independent identically distributed random variables. The main difficulty is that the distribution of $\Sbn_{1}$ depends on $n$. To simplify notation, set $f(t)= e^{- \frac{t^{2}}{2}}$ for $t \in \R$. By Fourier inversion we have, for every $x \in \mathbb{R}$,

\begin{equation}
 \label{eq:fourierb}  p(x)= \frac{1}{2\pi} \int_{-\infty}^{\infty} e^{-itx} f(t) {\d}t.
 \end{equation}
Also, for $k \in \mathbb{Z}$,
\[ \Pr{\Sbn_{N}=k} = \frac{1}{2\pi} \int_{-\pi}^{\pi} e^{-itk} \phibn(t)^{ N} {\mathrm{d}}t.\]
In the following, we only consider values of $x \in \R$ such that $N \frac{K_{n}+1}{n-K_{n}}+x  \DbnN$ is an integer.
For such $x$,
\begin{multline*}
  \DbnN \cdot \Pr{\Sbn_{ N}= N \frac{K_{n}+1}{n-K_{n}}+x  \DbnN}\\
  =  \frac{1}{2\pi} \int^{\pi \DbnN}_{-\pi \DbnN} e^{- i t x} \left(  \phibn \left(  \frac{t}{\DbnN} \right)   e^{-it \frac{K_{n}+1}{(n-K_{n}) \DbnN}} \right)^{N} {\mathrm{d}}t.
\end{multline*}
Therefore for every fixed $A>0$, for $n$ sufficiently large,
\begin{multline*}
  \left| \DbnN \cdot \Pr{\Sbn_{N}= N \frac{K_{n}+1}{n-K_{n}}+x  \DbnN} -p( x)\right| \\
  \leq \frac{1}{2 \pi} \big(|I^{(1)}_{A}(x,N,n)| + |I^{(2)}_{A}(x,N,n)|+|I^{(3)}(x)|\big),
\end{multline*}
where
\begin{align*}
  I^{(1)}_{A}(x,N,n)&=\int_{-A}^{A} e^{-i t x} \left(   \left(  \phibn \left(  \frac{t}{\DbnN} \right)   e^{-it \frac{K_{n}+1}{(n-K_{n}) \DbnN}} \right)^{N}- f(t) \right) {\mathrm{d}}t,\\
 I^{(2)}_{A}(x,N,n)&= \int_{A<|t|< \pi \DbnN} e^{-itx - it \frac{(K_{n}+1) N}{(n-K_{n}) \DbnN}} \phibn \left(  \frac{t}{\DbnN} \right) ^{ N}{\mathrm{d}}t, \\
 I^{(3)}_{A}(x)&= \int_{|t|>A} e^{-itx} f(t) {\mathrm{d}}t.
 \end{align*}
We shall check that for every fixed $\epsilon>0$, there exists $A>0$ such that for every $n$ sufficiently large, for every $ un \leq N \leq  n$, for every $x \in \R$ (with the above integrality condition),  $|I^{(1)}_{A}(x,N,n)| \leq \epsilon$, $|I^{(2)}_{A}(x,N,n)| \leq \epsilon$ and $|I^{(3)}_{A}(x)| \leq\epsilon$.



 We choose $A >0$ such that
\begin{equation}
\label{eq:Ab} 2 \int_{A}^{\infty } e^{-  \frac{1}{9} t^{2}} {\d}t < \epsilon.
\end{equation}


\emph{Bounding $|I^{(1)}_{A}(x,N,n)|$.} Since the convergence \eqref{eq:cv1b} holds  uniformly on compact subsets of $\R$, for   $n$ sufficiently large, for every $ un \leq N \leq n$ and $x \in \R$, we have $|I^{(1)}_{A}(x,N,n)| \leq \epsilon$.

\emph{Bounding $|I^{(3)}_{A}(x)|$.} By the choice of $A$ in \eqref{eq:Ab}
since $|f(t)|=e^{-  \frac{1}{2} t^{2}} \le e^{-  \frac{1}{9} t^{2}}$,
for every $x \in \R$ we have $|I^{(3)}_{A}(x)| \leq \epsilon$.


 \emph{Bounding $|I^{(2)}_{{A} }(x,N,n)|$.}  By Lemma~\ref{lem:phib},  for every $n$ sufficiently large, for every $ un \leq N \leq n$, $A<|t| \leq \pi  \DbnN$ we have
\[|I^{(2)}(x,N,n)| \leq 2 \int_{A}^{\pi \DbnN} e^{-  \frac{ K_{n} N }{ 8 n (\DbnN)^{2}} \cdot t^{2} } {\d}t \leq 2 \int_{A}^{\infty } e^{-  \frac{1}{9} t^{2}} {\d}t,\]
   which is less than $\epsilon$ by \eqref{eq:Ab}. This completes the proof.
\end{proof}





\paragraph*{Local limit theorem for $\Scn$.}

In the case of $\Scn$, the analysis is more subtle.
As above, we start with an estimate on the characteristic function
of the step distribution of the random walk
\[\phicn(t)\coloneqq\Es{e^{it\Scn_{1}}}=\frac{F( \bcn e^{it})}{F(\bcn)}.\]
To simplify notation, set  $y_{n}=1/e-\bcn$.

\begin{lemma}\label{lem:technique}
Assume that $\tfrac{K_{n}}{\sqrt{n}} \rightarrow c>0$ as $n \rightarrow \infty$.  There exist constants $A_{0},\epsilon,\kappa>0$, which may only depend on $c$, such that the following holds.
For every $n$ sufficiently large and for every $|t| \in (\frac{A_{0}}{\Dcn},\epsilon)$, we have $|\phicn(t)| \leq e^{- \kappa |t|^{1/2}}$.
\end{lemma}



\begin{proof}We start with some preliminary observations which fix the values of the different constants.
\begin{enumerate}
\item[--]  Since $\Dcn \sim n /c$ and $y_{n} \sim \frac{c^{2}}{2en}$ (as was seen in the proof of Lemma~\ref{lem:exist}),
  we can find a constant $\gamma>0$ depending only on $c$ such that
  the conditions $t \ge  {A_0}/{\Dcn}$ and $A_{0} \geq 1$ imply
  $t \ge {\gamma A_0}/{n}$ and $t \ge \gamma y_{n}$,
  for $n$ sufficiently large.
\item[--] We then fix $A_{0} \geq 1$ such that $\frac{2c}{\sqrt{\gamma A_{0}}} \leq  \frac{1}{4\sqrt{e}}$.
\item[--] From the expansion $ \ln F \big( \tfrac{1}{e}-z \big)= 1-\sqrt{2e} \sqrt{z}+o(\sqrt{z})$, there exists $\eta>0$ such that
\begin{equation}
\label{eq:z} \Re(z) \ge 0, \, |z| \leq \eta  \qquad \implies  \qquad  \Re \ln F \left( \frac{1}{e}-z \right) \leq  1-  \Re \sqrt{z}.
\end{equation}
Then set $\epsilon= \sqrt{\eta^{2}/(e^{-2}+\gamma^{-2})}$.
\end{enumerate}


We now turn to the main part of the proof. For $t \in \R$, since
$ \ln |\phicn(t)| =\Re \ln \phicn(t)$, we have
\begin{equation}
\label{eq:lnphi} \ln |\phicn(t)| =\Re \ln F( \bcn e^{it})- \ln F( \bcn).
\end{equation}
We first study the term $\ln F( \bcn)$.
By \eqref{eq:devF} and \eqref{eq:est_bcirc}, we have $\ln F( \bcn)=1- \tfrac{c}{\sqrt{n}}+ o \big(  \frac{1}{\sqrt{n}} \big)$.  As a consequence, for every $n$ sufficiently large, for every $  {A_{0}}/{\Dcn} \leq |t| \leq 1$, we have
\begin{equation}
\label{eq:ln}\ln F( \bcn) \geq 1- \frac{2c}{\sqrt{n}} \geq 1- \frac{2c}{\sqrt{\gamma A_{0}}} \sqrt{|t|}.
\end{equation}

Consider now the term $\Re \ln F( \bcn e^{it})$.
To this end, we set $z=z_{n}(t)=\tfrac{1}{e}-\bcn e^{it}$.
Note that we always have $\Re(z_n(t)) \ge 0$.
Observing that $e^{-it} z_{n}(t)= \tfrac{1}{e}(e^{-it}-1)+y_n$, we have
\begin{equation}
  \label{eq:module} |z_{n}(t)|^{2}=  \tfrac{1}{e^2} \big|e^{-it}-1\big|^2
  +\tfrac{2}{e} \big\langle y_n, e^{-it}-1 \big\rangle +|y_n|^2
  = \frac{2(1-\cos(t))}{e^{2}} -\frac{2 y_{n}(1-\cos(t))}{e}+y_{n}^{2}.
\end{equation}
For $n$ sufficiently large and $|t| \leq \epsilon$, it follows that
\begin{equation}
\label{ineq:module} |z_{n}(t)|^{2} \leq  \left( \frac{1}{e^{2}} + \frac{1}{\gamma^{2}} \right)  t^{2} \leq \eta^{2}.
\end{equation}
Using \eqref{eq:z}, we have
\[\Re \ln F( \bcn e^{it}) \leq 1-   \Re \sqrt{z_{n}(t)}.\]
Since $\Re z_{n}(t) \ge 0$, we have
\[ \Re \sqrt{z_{n}(t)}=\sqrt{\frac{\Re(z_n(t))+|z_n(t)|}{2}} \ge \sqrt{\frac{|z_n(t)|}{2}}. \]
But, by \eqref{eq:module}, for $n$ sufficiently large and  for every $|t|$ in $(A_{0}/\Dcn,1)$,
\[|z_{n}(t)|^{2} = \frac{2(1-\cos(t))}{e^{2}} (1-y_{n} e)+y_{n}^{2} \geq \frac{1-\cos(t)}{e^{2}} \geq  \frac{t^{2}}{4e^{2}}.\]
As a consequence $\Re \sqrt{z_{n}(t)} \geq \frac{\sqrt{|t|}}{2\sqrt{e}}$, so that
\begin{equation}
\label{eq:reln}\Re \ln F( \bcn e^{it}) \leq 1- \frac{\sqrt{|t|}}{2\sqrt{e}}.
\end{equation}
To conclude, by combining \eqref{eq:lnphi} with \eqref{eq:ln} and \eqref{eq:reln},
we get that for every $n$ sufficiently large, for every $|t|$ in $(A_{0}/\Dcn, \epsilon)$,
\[\ln|\phicn(t)| \leq \left(  \frac{2c}{\sqrt{\gamma A_{0}}}- \frac{1}{2\sqrt{e}} \right) \cdot \sqrt{|t|} \leq -\frac{1}{4\sqrt{e}} \sqrt{|t|}.\]
This completes the proof, with $\kappa= \frac{1}{4\sqrt{e}}$.
\end{proof}


We are now ready to tackle the proof of Lemma~\ref{lem:ll2}.
\begin{proof}[Proof of Lemma~\ref{lem:ll2}]
The first step is to show the convergence
\begin{equation}
\label{eq:cv1}\Es{e^{i t \frac{\Scn_{ u K_{n} +j}}{ \Dcn }}}  \quad \mathop{\longrightarrow}_{n \rightarrow \infty} \quad \exp \left( u c^{2} \left( 1 -\sqrt{1-  \frac{2it}{c}}  \right)  \right),
\end{equation}
uniformly for $t$ in compact subsets of $\R$ and $|j| \leq n^{3/8}$.

The characteristic function $\phicn(t)=\Es{e^{it\Scn_{1}}}$  of $\Scn_{1}$
is given as $\phicn(t)= \frac{F( \bcn e^{it})}{F(\bcn)}$,
where $\bcn$ has been chosen such that $ \frac{\bcn F'(\bcn)}{F(\bcn)}= \frac{n-K_{n}}{K_{n}+1}$.
Recall that $\bcn=1/e-y_{n}$ with
\hbox{$y_{n} \sim \frac{1}{2e}  \big( \frac{K_{n}}{n} \big)^{2} \sim \frac{c^{2}}{2e} \cdot \frac{1}{n}$}
(from \eqref{eq:est_bcirc}).
Since $ \Dcn \sim  n /c$, we have, uniformly on compact subsets of $\R$ as $n \rightarrow \infty$,
\[\bcn e^{ \frac{it}{\Dcn}}=(1/e-y_{n})e^{ \frac{it}{\Dcn}}=1/e-  \left( \frac{c^{2}}{2e} - \frac{it c}{e } \right) \cdot \frac{1}{{n}}+o \left(  \frac{1}{n} \right).\]
By \eqref{eq:devF}, we have $ \ln F \left( \frac{1}{e}-z \right)= 1-\sqrt{2e} \sqrt{z}+o(\sqrt{z})$ as $z \rightarrow 0$.
As a consequence as $n \rightarrow \infty$, \[\ln \phicn \left( \frac{t}{ \Dcn } \right)=-\sqrt{2e} \sqrt{ \frac{c^{2}}{2e} - \frac{it c}{e}} \cdot \frac{1}{\sqrt{n}} + \sqrt{2e} \sqrt{ \frac{c^{2}}{2e}} \cdot \frac{1}{\sqrt{n}}+o \left(  \frac{1}{\sqrt{n}} \right),\]
 uniformly on compact subsets of $\R$. Therefore
\[( uK_{n}  +j) \ln \phicn \left( \frac{t}{\Dcn} \right)  \quad \mathop{\longrightarrow}_{n \rightarrow \infty} \quad u c^{2}-uc \sqrt{c^{2}- 2itc},\]
 uniformly when $t$ belongs to compact subsets of $\R$ and $|j| \leq n^{3/8}$, which implies \eqref{eq:cv1}.

 As in the proof of Lemma~\ref{lem:ll1}, the second step consists in estimating $\Pr{\Scn_{ uK_{n}  +j}=k}$ by Fourier inversion.
 We first let, for $t \in \mathbb{R}$,
\[f(t)=\exp \left( u c^{2} \left( 1 -\sqrt{1-  \frac{2it}{c}}  \right)  \right)\]
be the expression appearing in \eqref{eq:cv1}.
By \cite[Sec.1.2.5]{Kyp06}),  $f$ is the characteristic function of a random variable with explicit density
\[q_u(x)=\left( \frac{u^{2} c^{3}}{2 \pi x^{3}} \right)^{1/2} \exp \left(  - \frac{c (x-uc)^{2} }{2x} \right) \mathbbm{1}_{x>0},\]
so that by Fourier inversion we have, for every $x \in \mathbb{R}$,
\begin{equation}
 \label{eq:fourier}
 q_u(x)=  \frac{1}{2\pi}\int_{-\infty}^{\infty}  {\mathrm{d}t} e^{-itx} \exp \left( u c^{2} \left( 1 -\sqrt{1-  \frac{2it}{c}}  \right)  \right).
 \end{equation}
On the other hand,
for $k \in \mathbb{Z}$,
\[ \Pr{\Scn_{ uK_{n}  +j}=k} = \frac{1}{2\pi} \int_{-\pi}^{\pi} e^{-itk} \phicn(t)^{ uK_{n}  +j} {\mathrm{d}}t.\]
Therefore, assuming that $x \in \R$ is chosen so that $x \Dcn$ is an integer, we get that
\[ \Dcn \cdot \Pr{\Scn_{ uK_{n}  +j}=x  \Dcn}=  \frac{1}{2\pi} \int^{\pi \Dcn}_{-\pi \Dcn} e^{- i t x} \phicn \left(  \frac{t}{\Dcn} \right) ^{ uK_{n}  +j} {\mathrm{d}}t.\]
We deduce that for every fixed $A>0$ and $0< \epsilon \leq 1$, for $n$ sufficiently large, it holds
\begin{multline*}
  \left| \Dcn \cdot \Pr{\Scn_{ uK_{n}  +j}=x  \Dcn} - q_{u}(x)\right| \\
  \leq \frac{1}{2 \pi} (|I^{(1)}_{A}(x,j,n)| + |I^{(2)}_{\epsilon,A}(x,j,n)|+|I^{(3)}_{\epsilon}(x,j,n)|+|I^{(4)}_{A}(x)|),
\end{multline*}
where
\[I^{(1)}_{A}(x,j,n)=\int_{-A}^{A} e^{-i t x} \left(  \phicn \left(  \frac{t}{\Dcn} \right) ^{ uK_{n}  +j}- f(t) \right) {\mathrm{d}}t\]
\[ I^{(2)}_{\epsilon,A}(x,j,n)= \int_{A<|t|< \epsilon \Dcn} e^{-itx} \phicn \left(  \frac{t}{\Dcn} \right) ^{ uK_{n}  +j}{\mathrm{d}}t,\]
\[ I^{(3)}_{\epsilon}(x,j,n)=\int_{\epsilon \Dcn<|t|<\pi \Dcn} e^{-itx} \phicn  \left(  \frac{t}{\Dcn} \right) ^{ uK_{n}  +j} {\mathrm{d}}t, \qquad I^{(4)}_{A}(x)= \int_{|t|>A} e^{-itx} f(t) {\mathrm{d}}t .\]
We shall check that for every fixed $\epsilon'>0$, there exists $A>0$ and $0 < \epsilon \leq 1$ such that for every $n$ sufficiently large, for every $|j| \leq n^{3/8}$, for every $x \in \R$,
 $|I^{(1)}_{A}(x,j,n)| \leq \epsilon', |I^{(2)}_{\epsilon,A}(x,j,n)| \leq \epsilon', |I^{(3)}_{\epsilon}(x,j,n)| \leq \epsilon', |I^{(4)}_{A}(x)| \leq\epsilon'$.


Fix $\epsilon'>0$. We first explain how to choose $A$ and $\epsilon$.
As a preliminary observation, note that there exists a constant $C_{1}>0$
(which only depend on $u$ and $c$) such that $|f(t)| \leq C_{1} e^{-C_{1}^{-1} \sqrt{t}}$ for $t \geq 0$
(this is straightforward, using the explicit expression of $f$).
Let $A_{0}$ and $\kappa$ be the constants given by Lemma~\ref{lem:technique}.
We may choose $A \geq A_{0}$ such that
\begin{equation}
\label{eq:A}2\int_{A}^{\infty} |f(t)|  \mathrm{d}t <\epsilon' \qquad \textrm{and} \qquad  2 \int_{A}^{\infty} e^{- \frac{\kappa u c^{3/2}}{2} t^{1/2}} \mathrm{d}t < \epsilon'.
\end{equation}
We then let $\epsilon$ be also given by Lemma~\ref{lem:technique}.






\emph{Bounding $|I^{(1)}_{A}(x,j,n)|$.} Since the convergence \eqref{eq:cv1} holds  uniformly on compact subsets of $\R$ and $|j| \leq n^{3/8}$, for   $n$ sufficiently large, for every $|j| \leq n^{3/8}$ and $x \in \R$, we have $|I^{(1)}_{A}(x,j,n)| \leq \epsilon'$.

\emph{Bounding $|I^{(4)}_{A}(x)|$.} By the choice of $A$ in \eqref{eq:A}, for every $n$ sufficiently large, for every $x \in \R$ we have $|I^{(4)}_{A}(x,n)| \leq \epsilon'$.

\emph{Bounding $|I^{(3)}_{\epsilon}(x,j,n)|$.}  Set $\phi(t)=\frac{F(e^{it}/e)}{F(1/e)}$, which is the characteristic function of  the probability distribution $ \frac{(i+1)^{(i-1)}}{e^{i+1} i!}$ on $\Z_{+}$. Since the latter is non lattice, we have $|\phi(t)|<1$ for $t \in (0,2\pi)$ (see e.g.~\cite[Theorem 3.5.1]{Dur10}).
Therefore, there exists $C_2>0$ such that $|\phi(t)|<e^{-2C_2}$ for $ \epsilon \leq |t| \leq \pi$.
Since $F$ is a power series with nonnegative coefficients and converges in $1/e$,
it converges uniformly on $\{z: |z| \le 1/e\}$ and is continuous on this compact set.
It is therefore uniformly continuous and
we have $\phicn(t) \rightarrow \phi(t)$, uniformly for $t \in \R$.
In particular, for $n$ sufficiently large and $t$ with $ \epsilon \leq |t| \leq \pi$,
we have $|\phicn(t)|<e^{-C_2}$.
 Therefore, for every $n$ sufficiently large, for every $|j| \leq n^{3/8}$ and  $x \in \R$,
 \[ \left|I^{(3)}_{\epsilon}(x,j,n) \right|  \leq 2 \int_{\epsilon \Dcn}^{\pi \Dcn} e^{-C_2 (uK_{n}  +n^{3/8})} \mathrm{d}t
 \leq 2 \pi \Dcn \cdot (e^{- C_2 u c \sqrt{n}/2} ) \le \epsilon'.\]

 \emph{Bounding $|I^{(2)}_{\epsilon, {A} }(x,j,n)|$.}  By Lemma~\ref{lem:technique}, for every $A<|t| \leq \epsilon \Dcn$ we have \[  \left| \phicn \left(  \frac{t}{\Dcn} \right) \right|  \leq e^{-\kappa \left(  \frac{|t|}{\Dcn} \right) ^{1/2}}.\]
Hence, for $n$ sufficiently large, for every $|j| \leq n^{3/8}$ and $x \in \R$,
 \[ \left| I^{(2)}_{\epsilon,A}(x,j,n)  \right| \leq 2 \int_{A}^{\epsilon \Dcn}\left( e^{-\kappa  \left(  \frac{t}{\Dcn} \right) ^{1/2}} \right) ^{ uK_{n}  -n^{3/8}} \mathrm{d}t .\]
 Since $K_{n} \sim c \sqrt{n}$ and $ \Dcn\sim  n /c$, it follows that for $n$ sufficiently large, for every $|j| \leq n^{3/8}$ and $x \in \R$,
  \[ \left| I^{(2)}_{\epsilon,A}(x,j,n) \right| \leq 2 \int_{A}^{\epsilon \Dcn}  e^{- \frac{\kappa u  c^{3/2}}{2} t^{1/2}} \mathrm{d}t  \leq 2 \int_{A}^{\infty} e^{- \frac{\kappa u c^{3/2}}{2} t^{1/2}} \mathrm{d}t,\]
  which is less than $\epsilon'$ by \eqref{eq:A}.   This completes the proof.
  \end{proof}
